La apertura a 25 agrupaciones permite mirar dentro de algunas de las grandes actividades económicas usadas en la sección anterior y repetir el mismo ejercicio de productividad laboral con una mayor desagregación. En cada actividad se cruza el PIB reportado por el DANE con ocupados y horas de la GEIH, de modo que el análisis mantiene las dos medidas centrales del informe: PIB por trabajador y PIB por hora trabajada.

\begin{table}[H]
\centering
\caption{Productividad laboral por actividad económica, 25 agrupaciones CIIU, 2010--2025}
\label{tab:pib_geih_productividad_25}
\footnotesize
\begin{tabular}{p{0.36\textwidth}rrrrrr}
\toprule
& \multicolumn{3}{c}{PIB por trabajador} & \multicolumn{3}{c}{PIB por hora} \\
& \multicolumn{3}{c}{\footnotesize Millones de pesos de 2015} & \multicolumn{3}{c}{\footnotesize Miles de pesos de 2015} \\
\cmidrule(lr){2-4}\cmidrule(lr){5-7}
Actividad económica & 2010 & 2025 & Crec. & 2010 & 2025 & Crec. \\
\midrule
Artes, entretenimiento y otros servicios & 13,8 & 32,8 & 5,95\% & 6,8 & 16,0 & 5,89\% \\
Información y comunicaciones & 45,3 & 78,6 & 3,74\% & 17,6 & 34,0 & 4,50\% \\
Agropecuario & 13,5 & 20,4 & 2,79\% & 5,8 & 9,4 & 3,27\% \\
Comercio y reparación & 15,2 & 22,8 & 2,75\% & 6,0 & 9,5 & 3,16\% \\
Salud y servicios sociales & 33,2 & 49,1 & 2,65\% & 14,5 & 21,3 & 2,58\% \\
Hogares como empleadores & 6,4 & 9,4 & 2,63\% & 2,7 & 4,5 & 3,49\% \\
Alimentos, bebidas y tabaco & 40,0 & 50,8 & 1,60\% & 16,9 & 21,8 & 1,70\% \\
Educación & 41,6 & 52,3 & 1,54\% & 21,4 & 25,8 & 1,26\% \\
Madera, papel e impresión & 37,7 & 46,9 & 1,46\% & 15,7 & 21,8 & 2,21\% \\
Financieras y seguros & 101,0 & 124,5 & 1,40\% & 43,0 & 54,2 & 1,56\% \\
Administración pública & 62,7 & 77,2 & 1,40\% & 23,7 & 32,2 & 2,06\% \\
Textiles, confecciones y cuero & 15,7 & 17,5 & 0,74\% & 6,5 & 7,8 & 1,19\% \\
Metalurgia, maquinaria y equipo & 45,1 & 48,8 & 0,53\% & 17,7 & 20,5 & 1,00\% \\
Inmobiliarias & 280,0 & 284,6 & 0,11\% & 105,3 & 119,5 & 0,85\% \\
Transporte y almacenamiento & 28,0 & 27,8 & -0,04\% & 9,4 & 10,9 & 0,94\% \\
Electricidad, gas y vapor & 267,5 & 262,2 & -0,13\% & 105,0 & 109,6 & 0,28\% \\
Refinación, químicos y minerales & 119,3 & 106,6 & -0,75\% & 46,8 & 44,8 & -0,30\% \\
Alojamiento y comida & 27,6 & 23,4 & -1,08\% & 11,4 & 10,7 & -0,42\% \\
Muebles y otras manufactureras & 21,2 & 17,4 & -1,31\% & 8,9 & 7,6 & -1,06\% \\
Profesionales y administrativas & 56,6 & 42,4 & -1,90\% & 26,0 & 20,6 & -1,52\% \\
Actividades especializadas de construcción & 42,1 & 26,3 & -3,10\% & 16,5 & 11,1 & -2,65\% \\
Obras civiles & 42,1 & 26,3 & -3,10\% & 16,5 & 11,1 & -2,65\% \\
Edificaciones & 42,1 & 26,3 & -3,10\% & 16,5 & 11,1 & -2,65\% \\
Minas & 200,8 & 115,8 & -3,60\% & 77,6 & 48,4 & -3,09\% \\
Agua, saneamiento y desechos & 106,3 & 37,9 & -6,64\% & 42,4 & 16,9 & -5,96\% \\
\bottomrule
\end{tabular}
\caption*{\footnotesize Nota: PIB por trabajador en millones de pesos constantes de 2015; PIB por hora en miles de pesos constantes de 2015. A nivel de actividad económica, el numerador corresponde estrictamente al valor agregado bruto de cada actividad. Ocupados y horas se agregan desde GEIH usando subramas homologadas. Cuando la apertura del DANE es más fina que la apertura laboral comparable, los ocupados y las horas se asignan proporcionalmente al PIB de cada subactividad dentro del grupo comparable. Fuente: cálculos propios con DANE y GEIH.}
\end{table}


\textbf{La desagregación a 25 agrupaciones muestra que la heterogeneidad de la productividad es mayor que la observada en las doce agrupaciones.} Las mayores tasas de crecimiento del PIB por trabajador se observan en Artes, entretenimiento y otros servicios (5,95\%); Información y comunicaciones (3,74\%); Agropecuario (2,79\%). En el otro extremo, las menores tasas aparecen en Agua, saneamiento y desechos (-6,64\%); Minas (-3,60\%); Edificaciones (-3,10\%). La lectura por hora confirma el mismo mensaje general: las mayores mejoras en PIB por hora se concentran en Artes, entretenimiento y otros servicios (5,89\%); Información y comunicaciones (4,50\%); Hogares como empleadores (3,49\%), mientras que los rezagos más marcados se ubican en Agua, saneamiento y desechos (-5,96\%); Minas (-3,09\%); Edificaciones (-2,65\%).

A continuación se presenta el detalle de cada agrupación. Las comparaciones se hacen frente al agregado nacional para mantener el mismo punto de referencia usado en la sección anterior.

\subsection{Agropecuario}

Esta agrupación reúne agricultura, ganadería, caza, silvicultura y pesca. En la CIIU Rev. 4 A.C. corresponde a A; divisiones 01--03.

\begin{table}[H]
\centering
\caption{Agropecuario: PIB, ocupados, horas y productividad laboral, 2010--2025}
\label{tab:agg25_a_productividad}
\small
\begin{tabular}{lrrr}
\toprule
Indicador & 2010 & 2025 & Crec. anual \\
\midrule
PIB real  (Billones de pesos de 2015) & 39,9 & 63,9 & 3,18\% \\
Ocupados (Millones) & 3,0 & 3,1 & 0,38\% \\
PIB por trabajador (Millones de pesos de 2015) & 13,5 & 20,4 & 2,79\% \\
Horas semanales por trabajador & 45,0 & 42,0 & -0,46\% \\
PIB por hora trabajada (Miles de pesos de 2015) & 5,8 & 9,4 & 3,27\% \\
\bottomrule
\end{tabular}
\end{table}

\textbf{Frente al agregado nacional, la comparación variable por variable muestra diferencias relevantes.} El PIB real de la actividad registró una tasa anualizada de 3,18\%, muy cerca del agregado (3,16\%). La ocupación en la actividad registró una tasa anualizada de 0,38\%, por debajo del agregado (2,09\%). El PIB por trabajador registró una tasa anualizada de 2,79\%, por encima del agregado (1,05\%). Las horas semanales por trabajador registraron una tasa anualizada de -0,46\%, muy cerca de la caída agregada (-0,45\%). El PIB por hora trabajada registró una tasa anualizada de 3,27\%, por encima del agregado (1,51\%).

\textbf{En niveles de 2025, el tamaño relativo y la productividad de la actividad económica también importan.} La actividad representó 6,3\% del PIB real agregado, con 63,9 billones de pesos de 2015, y concentró 13,6\% de los ocupados, con 3,1 millones de personas. El PIB por trabajador fue 20,4 millones de pesos de 2015 por ocupado, por debajo del agregado (44,5 millones). Las horas semanales por trabajador fueron 42,0, muy cerca del agregado (43,7). El PIB por hora trabajada fue 9,4 mil pesos de 2015 por hora, por debajo del agregado (19,6 mil).

\subsection{Minas}

Esta agrupación reúne explotación de minas y canteras. En la CIIU Rev. 4 A.C. corresponde a B; divisiones 05--09.

\begin{table}[H]
\centering
\caption{Minas: PIB, ocupados, horas y productividad laboral, 2010--2025}
\label{tab:agg25_b_productividad}
\small
\begin{tabular}{lrrr}
\toprule
Indicador & 2010 & 2025 & Crec. anual \\
\midrule
PIB real  (Billones de pesos de 2015) & 38,4 & 34,7 & -0,67\% \\
Ocupados (Millones) & 0,2 & 0,3 & 3,04\% \\
PIB por trabajador (Millones de pesos de 2015) & 200,8 & 115,8 & -3,60\% \\
Horas semanales por trabajador & 49,7 & 46,0 & -0,52\% \\
PIB por hora trabajada (Miles de pesos de 2015) & 77,6 & 48,4 & -3,09\% \\
\bottomrule
\end{tabular}
\end{table}

\textbf{Frente al agregado nacional, la comparación variable por variable muestra diferencias relevantes.} El PIB real de la actividad registró una tasa anualizada de -0,67\%, por debajo del agregado (3,16\%). La ocupación en la actividad registró una tasa anualizada de 3,04\%, por encima del agregado (2,09\%). El PIB por trabajador registró una tasa anualizada de -3,60\%, por debajo del agregado (1,05\%). Las horas semanales por trabajador registraron una tasa anualizada de -0,52\%, muy cerca de la caída agregada (-0,45\%). El PIB por hora trabajada registró una tasa anualizada de -3,09\%, por debajo del agregado (1,51\%).

\textbf{En niveles de 2025, el tamaño relativo y la productividad de la actividad económica también importan.} La actividad representó 3,4\% del PIB real agregado, con 34,7 billones de pesos de 2015, y concentró 1,3\% de los ocupados, con 0,3 millones de personas. El PIB por trabajador fue 115,8 millones de pesos de 2015 por ocupado, por encima del agregado (44,5 millones). Las horas semanales por trabajador fueron 46,0, por encima del agregado (43,7). El PIB por hora trabajada fue 48,4 mil pesos de 2015 por hora, por encima del agregado (19,6 mil).

\subsection{Alimentos, bebidas y tabaco}

Esta agrupación reúne elaboración de productos alimenticios, bebidas y tabaco. En la CIIU Rev. 4 A.C. corresponde a C; divisiones 10--12.

\begin{table}[H]
\centering
\caption{Alimentos, bebidas y tabaco: PIB, ocupados, horas y productividad laboral, 2010--2025}
\label{tab:agg25_c01_productividad}
\small
\begin{tabular}{lrrr}
\toprule
Indicador & 2010 & 2025 & Crec. anual \\
\midrule
PIB real  (Billones de pesos de 2015) & 23,9 & 33,5 & 2,26\% \\
Ocupados (Millones) & 0,6 & 0,7 & 0,65\% \\
PIB por trabajador (Millones de pesos de 2015) & 40,0 & 50,8 & 1,60\% \\
Horas semanales por trabajador & 45,4 & 44,8 & -0,10\% \\
PIB por hora trabajada (Miles de pesos de 2015) & 16,9 & 21,8 & 1,70\% \\
\bottomrule
\end{tabular}
\end{table}

\textbf{Frente al agregado nacional, la comparación variable por variable muestra diferencias relevantes.} El PIB real de la actividad registró una tasa anualizada de 2,26\%, por debajo del agregado (3,16\%). La ocupación en la actividad registró una tasa anualizada de 0,65\%, por debajo del agregado (2,09\%). El PIB por trabajador registró una tasa anualizada de 1,60\%, por encima del agregado (1,05\%). Las horas semanales por trabajador registraron una tasa anualizada de -0,10\%, una caída menos pronunciada que la del agregado (-0,45\%). El PIB por hora trabajada registró una tasa anualizada de 1,70\%, muy cerca del agregado (1,51\%).

\textbf{En niveles de 2025, el tamaño relativo y la productividad de la actividad económica también importan.} La actividad representó 3,3\% del PIB real agregado, con 33,5 billones de pesos de 2015, y concentró 2,9\% de los ocupados, con 0,7 millones de personas. El PIB por trabajador fue 50,8 millones de pesos de 2015 por ocupado, por encima del agregado (44,5 millones). Las horas semanales por trabajador fueron 44,8, muy cerca del agregado (43,7). El PIB por hora trabajada fue 21,8 mil pesos de 2015 por hora, por encima del agregado (19,6 mil).

\subsection{Textiles, confecciones y cuero}

Esta agrupación reúne textiles, confecciones, cuero, calzado y artículos de viaje. En la CIIU Rev. 4 A.C. corresponde a C; divisiones 13--15.

\begin{table}[H]
\centering
\caption{Textiles, confecciones y cuero: PIB, ocupados, horas y productividad laboral, 2010--2025}
\label{tab:agg25_c02_productividad}
\small
\begin{tabular}{lrrr}
\toprule
Indicador & 2010 & 2025 & Crec. anual \\
\midrule
PIB real  (Billones de pesos de 2015) & 9,6 & 11,9 & 1,42\% \\
Ocupados (Millones) & 0,6 & 0,7 & 0,68\% \\
PIB por trabajador (Millones de pesos de 2015) & 15,7 & 17,5 & 0,74\% \\
Horas semanales por trabajador & 46,1 & 43,1 & -0,45\% \\
PIB por hora trabajada (Miles de pesos de 2015) & 6,5 & 7,8 & 1,19\% \\
\bottomrule
\end{tabular}
\end{table}

\textbf{Frente al agregado nacional, la comparación variable por variable muestra diferencias relevantes.} El PIB real de la actividad registró una tasa anualizada de 1,42\%, por debajo del agregado (3,16\%). La ocupación en la actividad registró una tasa anualizada de 0,68\%, por debajo del agregado (2,09\%). El PIB por trabajador registró una tasa anualizada de 0,74\%, por debajo del agregado (1,05\%). Las horas semanales por trabajador registraron una tasa anualizada de -0,45\%, muy cerca de la caída agregada (-0,45\%). El PIB por hora trabajada registró una tasa anualizada de 1,19\%, por debajo del agregado (1,51\%).

\textbf{En niveles de 2025, el tamaño relativo y la productividad de la actividad económica también importan.} La actividad representó 1,2\% del PIB real agregado, con 11,9 billones de pesos de 2015, y concentró 3,0\% de los ocupados, con 0,7 millones de personas. El PIB por trabajador fue 17,5 millones de pesos de 2015 por ocupado, por debajo del agregado (44,5 millones). Las horas semanales por trabajador fueron 43,1, muy cerca del agregado (43,7). El PIB por hora trabajada fue 7,8 mil pesos de 2015 por hora, por debajo del agregado (19,6 mil).

\subsection{Madera, papel e impresión}

Esta agrupación reúne madera, papel, cartón, impresión y reproducción de grabaciones. En la CIIU Rev. 4 A.C. corresponde a C; divisiones 16--18.

\begin{table}[H]
\centering
\caption{Madera, papel e impresión: PIB, ocupados, horas y productividad laboral, 2010--2025}
\label{tab:agg25_c03_productividad}
\small
\begin{tabular}{lrrr}
\toprule
Indicador & 2010 & 2025 & Crec. anual \\
\midrule
PIB real  (Billones de pesos de 2015) & 5,4 & 6,4 & 1,17\% \\
Ocupados (Millones) & 0,1 & 0,1 & -0,29\% \\
PIB por trabajador (Millones de pesos de 2015) & 37,7 & 46,9 & 1,46\% \\
Horas semanales por trabajador & 46,1 & 41,3 & -0,73\% \\
PIB por hora trabajada (Miles de pesos de 2015) & 15,7 & 21,8 & 2,21\% \\
\bottomrule
\end{tabular}
\end{table}

\textbf{Frente al agregado nacional, la comparación variable por variable muestra diferencias relevantes.} El PIB real de la actividad registró una tasa anualizada de 1,17\%, por debajo del agregado (3,16\%). La ocupación en la actividad registró una tasa anualizada de -0,29\%, por debajo del agregado (2,09\%). El PIB por trabajador registró una tasa anualizada de 1,46\%, por encima del agregado (1,05\%). Las horas semanales por trabajador registraron una tasa anualizada de -0,73\%, una caída más pronunciada que la del agregado (-0,45\%). El PIB por hora trabajada registró una tasa anualizada de 2,21\%, por encima del agregado (1,51\%).

\textbf{En niveles de 2025, el tamaño relativo y la productividad de la actividad económica también importan.} La actividad representó 0,6\% del PIB real agregado, con 6,4 billones de pesos de 2015, y concentró 0,6\% de los ocupados, con 0,1 millones de personas. El PIB por trabajador fue 46,9 millones de pesos de 2015 por ocupado, por encima del agregado (44,5 millones). Las horas semanales por trabajador fueron 41,3, por debajo del agregado (43,7). El PIB por hora trabajada fue 21,8 mil pesos de 2015 por hora, por encima del agregado (19,6 mil).

\subsection{Refinación, químicos y minerales}

Esta agrupación reúne refinación, químicos, farmacéuticos, caucho, plástico y minerales no metálicos. En la CIIU Rev. 4 A.C. corresponde a C; divisiones 19--23.

\begin{table}[H]
\centering
\caption{Refinación, químicos y minerales: PIB, ocupados, horas y productividad laboral, 2010--2025}
\label{tab:agg25_c04_productividad}
\small
\begin{tabular}{lrrr}
\toprule
Indicador & 2010 & 2025 & Crec. anual \\
\midrule
PIB real  (Billones de pesos de 2015) & 30,8 & 39,6 & 1,69\% \\
Ocupados (Millones) & 0,3 & 0,4 & 2,46\% \\
PIB por trabajador (Millones de pesos de 2015) & 119,3 & 106,6 & -0,75\% \\
Horas semanales por trabajador & 49,0 & 45,8 & -0,45\% \\
PIB por hora trabajada (Miles de pesos de 2015) & 46,8 & 44,8 & -0,30\% \\
\bottomrule
\end{tabular}
\end{table}

\textbf{Frente al agregado nacional, la comparación variable por variable muestra diferencias relevantes.} El PIB real de la actividad registró una tasa anualizada de 1,69\%, por debajo del agregado (3,16\%). La ocupación en la actividad registró una tasa anualizada de 2,46\%, por encima del agregado (2,09\%). El PIB por trabajador registró una tasa anualizada de -0,75\%, por debajo del agregado (1,05\%). Las horas semanales por trabajador registraron una tasa anualizada de -0,45\%, muy cerca de la caída agregada (-0,45\%). El PIB por hora trabajada registró una tasa anualizada de -0,30\%, por debajo del agregado (1,51\%).

\textbf{En niveles de 2025, el tamaño relativo y la productividad de la actividad económica también importan.} La actividad representó 3,9\% del PIB real agregado, con 39,6 billones de pesos de 2015, y concentró 1,6\% de los ocupados, con 0,4 millones de personas. El PIB por trabajador fue 106,6 millones de pesos de 2015 por ocupado, por encima del agregado (44,5 millones). Las horas semanales por trabajador fueron 45,8, muy cerca del agregado (43,7). El PIB por hora trabajada fue 44,8 mil pesos de 2015 por hora, por encima del agregado (19,6 mil).

\subsection{Metalurgia, maquinaria y equipo}

Esta agrupación reúne metalurgia, productos metálicos, equipo eléctrico, productos electrónicos, maquinaria, vehículos y equipo de transporte. En la CIIU Rev. 4 A.C. corresponde a C; divisiones 24--30 y 33.

\begin{table}[H]
\centering
\caption{Metalurgia, maquinaria y equipo: PIB, ocupados, horas y productividad laboral, 2010--2025}
\label{tab:agg25_c05_productividad}
\small
\begin{tabular}{lrrr}
\toprule
Indicador & 2010 & 2025 & Crec. anual \\
\midrule
PIB real  (Billones de pesos de 2015) & 13,1 & 15,5 & 1,10\% \\
Ocupados (Millones) & 0,3 & 0,3 & 0,57\% \\
PIB por trabajador (Millones de pesos de 2015) & 45,1 & 48,8 & 0,53\% \\
Horas semanales por trabajador & 49,1 & 45,7 & -0,47\% \\
PIB por hora trabajada (Miles de pesos de 2015) & 17,7 & 20,5 & 1,00\% \\
\bottomrule
\end{tabular}
\end{table}

\textbf{Frente al agregado nacional, la comparación variable por variable muestra diferencias relevantes.} El PIB real de la actividad registró una tasa anualizada de 1,10\%, por debajo del agregado (3,16\%). La ocupación en la actividad registró una tasa anualizada de 0,57\%, por debajo del agregado (2,09\%). El PIB por trabajador registró una tasa anualizada de 0,53\%, por debajo del agregado (1,05\%). Las horas semanales por trabajador registraron una tasa anualizada de -0,47\%, muy cerca de la caída agregada (-0,45\%). El PIB por hora trabajada registró una tasa anualizada de 1,00\%, por debajo del agregado (1,51\%).

\textbf{En niveles de 2025, el tamaño relativo y la productividad de la actividad económica también importan.} La actividad representó 1,5\% del PIB real agregado, con 15,5 billones de pesos de 2015, y concentró 1,4\% de los ocupados, con 0,3 millones de personas. El PIB por trabajador fue 48,8 millones de pesos de 2015 por ocupado, por encima del agregado (44,5 millones). Las horas semanales por trabajador fueron 45,7, muy cerca del agregado (43,7). El PIB por hora trabajada fue 20,5 mil pesos de 2015 por hora, muy cerca del agregado (19,6 mil).

\subsection{Muebles y otras manufactureras}

Esta agrupación reúne muebles, colchones, somieres y otras industrias manufactureras. En la CIIU Rev. 4 A.C. corresponde a C; divisiones 31--32.

\begin{table}[H]
\centering
\caption{Muebles y otras manufactureras: PIB, ocupados, horas y productividad laboral, 2010--2025}
\label{tab:agg25_c06_productividad}
\small
\begin{tabular}{lrrr}
\toprule
Indicador & 2010 & 2025 & Crec. anual \\
\midrule
PIB real  (Billones de pesos de 2015) & 5,0 & 6,0 & 1,28\% \\
Ocupados (Millones) & 0,2 & 0,3 & 2,63\% \\
PIB por trabajador (Millones de pesos de 2015) & 21,2 & 17,4 & -1,31\% \\
Horas semanales por trabajador & 45,9 & 44,1 & -0,25\% \\
PIB por hora trabajada (Miles de pesos de 2015) & 8,9 & 7,6 & -1,06\% \\
\bottomrule
\end{tabular}
\end{table}

\textbf{Frente al agregado nacional, la comparación variable por variable muestra diferencias relevantes.} El PIB real de la actividad registró una tasa anualizada de 1,28\%, por debajo del agregado (3,16\%). La ocupación en la actividad registró una tasa anualizada de 2,63\%, por encima del agregado (2,09\%). El PIB por trabajador registró una tasa anualizada de -1,31\%, por debajo del agregado (1,05\%). Las horas semanales por trabajador registraron una tasa anualizada de -0,25\%, muy cerca de la caída agregada (-0,45\%). El PIB por hora trabajada registró una tasa anualizada de -1,06\%, por debajo del agregado (1,51\%).

\textbf{En niveles de 2025, el tamaño relativo y la productividad de la actividad económica también importan.} La actividad representó 0,6\% del PIB real agregado, con 6,0 billones de pesos de 2015, y concentró 1,5\% de los ocupados, con 0,3 millones de personas. El PIB por trabajador fue 17,4 millones de pesos de 2015 por ocupado, por debajo del agregado (44,5 millones). Las horas semanales por trabajador fueron 44,1, muy cerca del agregado (43,7). El PIB por hora trabajada fue 7,6 mil pesos de 2015 por hora, por debajo del agregado (19,6 mil).

\subsection{Electricidad, gas y vapor}

Esta agrupación reúne suministro de electricidad, gas, vapor y aire acondicionado. En la CIIU Rev. 4 A.C. corresponde a D; división 35.

\begin{table}[H]
\centering
\caption{Electricidad, gas y vapor: PIB, ocupados, horas y productividad laboral, 2010--2025}
\label{tab:agg25_d_productividad}
\small
\begin{tabular}{lrrr}
\toprule
Indicador & 2010 & 2025 & Crec. anual \\
\midrule
PIB real  (Billones de pesos de 2015) & 14,9 & 21,9 & 2,60\% \\
Ocupados (Millones) & 0,1 & 0,1 & 2,73\% \\
PIB por trabajador (Millones de pesos de 2015) & 267,5 & 262,2 & -0,13\% \\
Horas semanales por trabajador & 49,0 & 46,0 & -0,41\% \\
PIB por hora trabajada (Miles de pesos de 2015) & 105,0 & 109,6 & 0,28\% \\
\bottomrule
\end{tabular}
\end{table}

\textbf{Frente al agregado nacional, la comparación variable por variable muestra diferencias relevantes.} El PIB real de la actividad registró una tasa anualizada de 2,60\%, por debajo del agregado (3,16\%). La ocupación en la actividad registró una tasa anualizada de 2,73\%, por encima del agregado (2,09\%). El PIB por trabajador registró una tasa anualizada de -0,13\%, por debajo del agregado (1,05\%). Las horas semanales por trabajador registraron una tasa anualizada de -0,41\%, muy cerca de la caída agregada (-0,45\%). El PIB por hora trabajada registró una tasa anualizada de 0,28\%, por debajo del agregado (1,51\%).

\textbf{En niveles de 2025, el tamaño relativo y la productividad de la actividad económica también importan.} La actividad representó 2,1\% del PIB real agregado, con 21,9 billones de pesos de 2015, y concentró 0,4\% de los ocupados, con 0,1 millones de personas. El PIB por trabajador fue 262,2 millones de pesos de 2015 por ocupado, por encima del agregado (44,5 millones). Las horas semanales por trabajador fueron 46,0, por encima del agregado (43,7). El PIB por hora trabajada fue 109,6 mil pesos de 2015 por hora, por encima del agregado (19,6 mil).

\subsection{Agua, saneamiento y desechos}

Esta agrupación reúne agua, saneamiento, residuos, desechos y saneamiento ambiental. En la CIIU Rev. 4 A.C. corresponde a E; divisiones 36--39.

\begin{table}[H]
\centering
\caption{Agua, saneamiento y desechos: PIB, ocupados, horas y productividad laboral, 2010--2025}
\label{tab:agg25_e_productividad}
\small
\begin{tabular}{lrrr}
\toprule
Indicador & 2010 & 2025 & Crec. anual \\
\midrule
PIB real  (Billones de pesos de 2015) & 7,0 & 8,4 & 1,19\% \\
Ocupados (Millones) & 0,1 & 0,2 & 8,38\% \\
PIB por trabajador (Millones de pesos de 2015) & 106,3 & 37,9 & -6,64\% \\
Horas semanales por trabajador & 48,2 & 43,3 & -0,72\% \\
PIB por hora trabajada (Miles de pesos de 2015) & 42,4 & 16,9 & -5,96\% \\
\bottomrule
\end{tabular}
\end{table}

\textbf{Frente al agregado nacional, la comparación variable por variable muestra diferencias relevantes.} El PIB real de la actividad registró una tasa anualizada de 1,19\%, por debajo del agregado (3,16\%). La ocupación en la actividad registró una tasa anualizada de 8,38\%, por encima del agregado (2,09\%). El PIB por trabajador registró una tasa anualizada de -6,64\%, por debajo del agregado (1,05\%). Las horas semanales por trabajador registraron una tasa anualizada de -0,72\%, una caída más pronunciada que la del agregado (-0,45\%). El PIB por hora trabajada registró una tasa anualizada de -5,96\%, por debajo del agregado (1,51\%).

\textbf{En niveles de 2025, el tamaño relativo y la productividad de la actividad económica también importan.} La actividad representó 0,8\% del PIB real agregado, con 8,4 billones de pesos de 2015, y concentró 1,0\% de los ocupados, con 0,2 millones de personas. El PIB por trabajador fue 37,9 millones de pesos de 2015 por ocupado, por debajo del agregado (44,5 millones). Las horas semanales por trabajador fueron 43,3, muy cerca del agregado (43,7). El PIB por hora trabajada fue 16,9 mil pesos de 2015 por hora, por debajo del agregado (19,6 mil).

\subsection{Edificaciones}

Esta agrupación reúne construcción de edificaciones residenciales y no residenciales. En la CIIU Rev. 4 A.C. corresponde a F; división 41. Para esta apertura de construcción, los ocupados y las horas se asignan desde la agrupación laboral de construcción según la participación de cada subactividad en el PIB de construcción de cada año.

\begin{table}[H]
\centering
\caption{Edificaciones: PIB, ocupados, horas y productividad laboral, 2010--2025}
\label{tab:agg25_f01_productividad}
\small
\begin{tabular}{lrrr}
\toprule
Indicador & 2010 & 2025 & Crec. anual \\
\midrule
PIB real  (Billones de pesos de 2015) & 22,0 & 20,5 & -0,49\% \\
Ocupados (Millones) & 0,5 & 0,8 & 2,69\% \\
PIB por trabajador (Millones de pesos de 2015) & 42,1 & 26,3 & -3,10\% \\
Horas semanales por trabajador & 48,9 & 45,7 & -0,46\% \\
PIB por hora trabajada (Miles de pesos de 2015) & 16,5 & 11,1 & -2,65\% \\
\bottomrule
\end{tabular}
\end{table}

\textbf{Frente al agregado nacional, la comparación variable por variable muestra diferencias relevantes.} El PIB real de la actividad registró una tasa anualizada de -0,49\%, por debajo del agregado (3,16\%). La ocupación en la actividad registró una tasa anualizada de 2,69\%, por encima del agregado (2,09\%). El PIB por trabajador registró una tasa anualizada de -3,10\%, por debajo del agregado (1,05\%). Las horas semanales por trabajador registraron una tasa anualizada de -0,46\%, muy cerca de la caída agregada (-0,45\%). El PIB por hora trabajada registró una tasa anualizada de -2,65\%, por debajo del agregado (1,51\%).

\textbf{En niveles de 2025, el tamaño relativo y la productividad de la actividad económica también importan.} La actividad representó 2,0\% del PIB real agregado, con 20,5 billones de pesos de 2015, y concentró 3,4\% de los ocupados, con 0,8 millones de personas. El PIB por trabajador fue 26,3 millones de pesos de 2015 por ocupado, por debajo del agregado (44,5 millones). Las horas semanales por trabajador fueron 45,7, muy cerca del agregado (43,7). El PIB por hora trabajada fue 11,1 mil pesos de 2015 por hora, por debajo del agregado (19,6 mil).

\subsection{Obras civiles}

Esta agrupación reúne construcción de carreteras, vías férreas, proyectos de servicio público y obras de ingeniería civil. En la CIIU Rev. 4 A.C. corresponde a F; división 42. Para esta apertura de construcción, los ocupados y las horas se asignan desde la agrupación laboral de construcción según la participación de cada subactividad en el PIB de construcción de cada año.

\begin{table}[H]
\centering
\caption{Obras civiles: PIB, ocupados, horas y productividad laboral, 2010--2025}
\label{tab:agg25_f02_productividad}
\small
\begin{tabular}{lrrr}
\toprule
Indicador & 2010 & 2025 & Crec. anual \\
\midrule
PIB real  (Billones de pesos de 2015) & 9,8 & 11,9 & 1,36\% \\
Ocupados (Millones) & 0,2 & 0,5 & 4,60\% \\
PIB por trabajador (Millones de pesos de 2015) & 42,1 & 26,3 & -3,10\% \\
Horas semanales por trabajador & 48,9 & 45,7 & -0,46\% \\
PIB por hora trabajada (Miles de pesos de 2015) & 16,5 & 11,1 & -2,65\% \\
\bottomrule
\end{tabular}
\end{table}

\textbf{Frente al agregado nacional, la comparación variable por variable muestra diferencias relevantes.} El PIB real de la actividad registró una tasa anualizada de 1,36\%, por debajo del agregado (3,16\%). La ocupación en la actividad registró una tasa anualizada de 4,60\%, por encima del agregado (2,09\%). El PIB por trabajador registró una tasa anualizada de -3,10\%, por debajo del agregado (1,05\%). Las horas semanales por trabajador registraron una tasa anualizada de -0,46\%, muy cerca de la caída agregada (-0,45\%). El PIB por hora trabajada registró una tasa anualizada de -2,65\%, por debajo del agregado (1,51\%).

\textbf{En niveles de 2025, el tamaño relativo y la productividad de la actividad económica también importan.} La actividad representó 1,2\% del PIB real agregado, con 11,9 billones de pesos de 2015, y concentró 2,0\% de los ocupados, con 0,5 millones de personas. El PIB por trabajador fue 26,3 millones de pesos de 2015 por ocupado, por debajo del agregado (44,5 millones). Las horas semanales por trabajador fueron 45,7, muy cerca del agregado (43,7). El PIB por hora trabajada fue 11,1 mil pesos de 2015 por hora, por debajo del agregado (19,6 mil).

\subsection{Actividades especializadas de construcción}

Esta agrupación reúne actividades especializadas para construcción de edificaciones y obras de ingeniería civil. En la CIIU Rev. 4 A.C. corresponde a F; división 43. Para esta apertura de construcción, los ocupados y las horas se asignan desde la agrupación laboral de construcción según la participación de cada subactividad en el PIB de construcción de cada año.

\begin{table}[H]
\centering
\caption{Actividades especializadas de construcción: PIB, ocupados, horas y productividad laboral, 2010--2025}
\label{tab:agg25_f03_productividad}
\small
\begin{tabular}{lrrr}
\toprule
Indicador & 2010 & 2025 & Crec. anual \\
\midrule
PIB real  (Billones de pesos de 2015) & 8,5 & 9,1 & 0,41\% \\
Ocupados (Millones) & 0,2 & 0,3 & 3,62\% \\
PIB por trabajador (Millones de pesos de 2015) & 42,1 & 26,3 & -3,10\% \\
Horas semanales por trabajador & 48,9 & 45,7 & -0,46\% \\
PIB por hora trabajada (Miles de pesos de 2015) & 16,5 & 11,1 & -2,65\% \\
\bottomrule
\end{tabular}
\end{table}

\textbf{Frente al agregado nacional, la comparación variable por variable muestra diferencias relevantes.} El PIB real de la actividad registró una tasa anualizada de 0,41\%, por debajo del agregado (3,16\%). La ocupación en la actividad registró una tasa anualizada de 3,62\%, por encima del agregado (2,09\%). El PIB por trabajador registró una tasa anualizada de -3,10\%, por debajo del agregado (1,05\%). Las horas semanales por trabajador registraron una tasa anualizada de -0,46\%, muy cerca de la caída agregada (-0,45\%). El PIB por hora trabajada registró una tasa anualizada de -2,65\%, por debajo del agregado (1,51\%).

\textbf{En niveles de 2025, el tamaño relativo y la productividad de la actividad económica también importan.} La actividad representó 0,9\% del PIB real agregado, con 9,1 billones de pesos de 2015, y concentró 1,5\% de los ocupados, con 0,3 millones de personas. El PIB por trabajador fue 26,3 millones de pesos de 2015 por ocupado, por debajo del agregado (44,5 millones). Las horas semanales por trabajador fueron 45,7, muy cerca del agregado (43,7). El PIB por hora trabajada fue 11,1 mil pesos de 2015 por hora, por debajo del agregado (19,6 mil).

\subsection{Comercio y reparación}

Esta agrupación reúne comercio al por mayor y al por menor, y reparación de vehículos automotores y motocicletas. En la CIIU Rev. 4 A.C. corresponde a G; divisiones 45--47.

\begin{table}[H]
\centering
\caption{Comercio y reparación: PIB, ocupados, horas y productividad laboral, 2010--2025}
\label{tab:agg25_g_productividad}
\small
\begin{tabular}{lrrr}
\toprule
Indicador & 2010 & 2025 & Crec. anual \\
\midrule
PIB real  (Billones de pesos de 2015) & 51,2 & 89,0 & 3,76\% \\
Ocupados (Millones) & 3,4 & 3,9 & 0,99\% \\
PIB por trabajador (Millones de pesos de 2015) & 15,2 & 22,8 & 2,75\% \\
Horas semanales por trabajador & 48,9 & 46,0 & -0,40\% \\
PIB por hora trabajada (Miles de pesos de 2015) & 6,0 & 9,5 & 3,16\% \\
\bottomrule
\end{tabular}
\end{table}

\textbf{Frente al agregado nacional, la comparación variable por variable muestra diferencias relevantes.} El PIB real de la actividad registró una tasa anualizada de 3,76\%, por encima del agregado (3,16\%). La ocupación en la actividad registró una tasa anualizada de 0,99\%, por debajo del agregado (2,09\%). El PIB por trabajador registró una tasa anualizada de 2,75\%, por encima del agregado (1,05\%). Las horas semanales por trabajador registraron una tasa anualizada de -0,40\%, muy cerca de la caída agregada (-0,45\%). El PIB por hora trabajada registró una tasa anualizada de 3,16\%, por encima del agregado (1,51\%).

\textbf{En niveles de 2025, el tamaño relativo y la productividad de la actividad económica también importan.} La actividad representó 8,7\% del PIB real agregado, con 89,0 billones de pesos de 2015, y concentró 17,0\% de los ocupados, con 3,9 millones de personas. El PIB por trabajador fue 22,8 millones de pesos de 2015 por ocupado, por debajo del agregado (44,5 millones). Las horas semanales por trabajador fueron 46,0, por encima del agregado (43,7). El PIB por hora trabajada fue 9,5 mil pesos de 2015 por hora, por debajo del agregado (19,6 mil).

\subsection{Transporte y almacenamiento}

Esta agrupación reúne transporte y almacenamiento. En la CIIU Rev. 4 A.C. corresponde a H; divisiones 49--53.

\begin{table}[H]
\centering
\caption{Transporte y almacenamiento: PIB, ocupados, horas y productividad laboral, 2010--2025}
\label{tab:agg25_h_productividad}
\small
\begin{tabular}{lrrr}
\toprule
Indicador & 2010 & 2025 & Crec. anual \\
\midrule
PIB real  (Billones de pesos de 2015) & 32,8 & 49,8 & 2,81\% \\
Ocupados (Millones) & 1,2 & 1,8 & 2,86\% \\
PIB por trabajador (Millones de pesos de 2015) & 28,0 & 27,8 & -0,04\% \\
Horas semanales por trabajador & 56,9 & 49,1 & -0,98\% \\
PIB por hora trabajada (Miles de pesos de 2015) & 9,4 & 10,9 & 0,94\% \\
\bottomrule
\end{tabular}
\end{table}

\textbf{Frente al agregado nacional, la comparación variable por variable muestra diferencias relevantes.} El PIB real de la actividad registró una tasa anualizada de 2,81\%, por debajo del agregado (3,16\%). La ocupación en la actividad registró una tasa anualizada de 2,86\%, por encima del agregado (2,09\%). El PIB por trabajador registró una tasa anualizada de -0,04\%, por debajo del agregado (1,05\%). Las horas semanales por trabajador registraron una tasa anualizada de -0,98\%, una caída más pronunciada que la del agregado (-0,45\%). El PIB por hora trabajada registró una tasa anualizada de 0,94\%, por debajo del agregado (1,51\%).

\textbf{En niveles de 2025, el tamaño relativo y la productividad de la actividad económica también importan.} La actividad representó 4,9\% del PIB real agregado, con 49,8 billones de pesos de 2015, y concentró 7,8\% de los ocupados, con 1,8 millones de personas. El PIB por trabajador fue 27,8 millones de pesos de 2015 por ocupado, por debajo del agregado (44,5 millones). Las horas semanales por trabajador fueron 49,1, por encima del agregado (43,7). El PIB por hora trabajada fue 10,9 mil pesos de 2015 por hora, por debajo del agregado (19,6 mil).

\subsection{Alojamiento y comida}

Esta agrupación reúne alojamiento y servicios de comida. En la CIIU Rev. 4 A.C. corresponde a I; divisiones 55--56.

\begin{table}[H]
\centering
\caption{Alojamiento y comida: PIB, ocupados, horas y productividad laboral, 2010--2025}
\label{tab:agg25_i_productividad}
\small
\begin{tabular}{lrrr}
\toprule
Indicador & 2010 & 2025 & Crec. anual \\
\midrule
PIB real  (Billones de pesos de 2015) & 23,8 & 40,6 & 3,63\% \\
Ocupados (Millones) & 0,9 & 1,7 & 4,76\% \\
PIB por trabajador (Millones de pesos de 2015) & 27,6 & 23,4 & -1,08\% \\
Horas semanales por trabajador & 46,4 & 42,0 & -0,66\% \\
PIB por hora trabajada (Miles de pesos de 2015) & 11,4 & 10,7 & -0,42\% \\
\bottomrule
\end{tabular}
\end{table}

\textbf{Frente al agregado nacional, la comparación variable por variable muestra diferencias relevantes.} El PIB real de la actividad registró una tasa anualizada de 3,63\%, por encima del agregado (3,16\%). La ocupación en la actividad registró una tasa anualizada de 4,76\%, por encima del agregado (2,09\%). El PIB por trabajador registró una tasa anualizada de -1,08\%, por debajo del agregado (1,05\%). Las horas semanales por trabajador registraron una tasa anualizada de -0,66\%, una caída más pronunciada que la del agregado (-0,45\%). El PIB por hora trabajada registró una tasa anualizada de -0,42\%, por debajo del agregado (1,51\%).

\textbf{En niveles de 2025, el tamaño relativo y la productividad de la actividad económica también importan.} La actividad representó 4,0\% del PIB real agregado, con 40,6 billones de pesos de 2015, y concentró 7,6\% de los ocupados, con 1,7 millones de personas. El PIB por trabajador fue 23,4 millones de pesos de 2015 por ocupado, por debajo del agregado (44,5 millones). Las horas semanales por trabajador fueron 42,0, muy cerca del agregado (43,7). El PIB por hora trabajada fue 10,7 mil pesos de 2015 por hora, por debajo del agregado (19,6 mil).

\subsection{Información y comunicaciones}

Esta agrupación reúne información y comunicaciones. En la CIIU Rev. 4 A.C. corresponde a J; divisiones 58--63.

\begin{table}[H]
\centering
\caption{Información y comunicaciones: PIB, ocupados, horas y productividad laboral, 2010--2025}
\label{tab:agg25_j_productividad}
\small
\begin{tabular}{lrrr}
\toprule
Indicador & 2010 & 2025 & Crec. anual \\
\midrule
PIB real  (Billones de pesos de 2015) & 18,3 & 31,3 & 3,66\% \\
Ocupados (Millones) & 0,4 & 0,4 & -0,08\% \\
PIB por trabajador (Millones de pesos de 2015) & 45,3 & 78,6 & 3,74\% \\
Horas semanales por trabajador & 49,6 & 44,4 & -0,73\% \\
PIB por hora trabajada (Miles de pesos de 2015) & 17,6 & 34,0 & 4,50\% \\
\bottomrule
\end{tabular}
\end{table}

\textbf{Frente al agregado nacional, la comparación variable por variable muestra diferencias relevantes.} El PIB real de la actividad registró una tasa anualizada de 3,66\%, por encima del agregado (3,16\%). La ocupación en la actividad registró una tasa anualizada de -0,08\%, por debajo del agregado (2,09\%). El PIB por trabajador registró una tasa anualizada de 3,74\%, por encima del agregado (1,05\%). Las horas semanales por trabajador registraron una tasa anualizada de -0,73\%, una caída más pronunciada que la del agregado (-0,45\%). El PIB por hora trabajada registró una tasa anualizada de 4,50\%, por encima del agregado (1,51\%).

\textbf{En niveles de 2025, el tamaño relativo y la productividad de la actividad económica también importan.} La actividad representó 3,1\% del PIB real agregado, con 31,3 billones de pesos de 2015, y concentró 1,7\% de los ocupados, con 0,4 millones de personas. El PIB por trabajador fue 78,6 millones de pesos de 2015 por ocupado, por encima del agregado (44,5 millones). Las horas semanales por trabajador fueron 44,4, muy cerca del agregado (43,7). El PIB por hora trabajada fue 34,0 mil pesos de 2015 por hora, por encima del agregado (19,6 mil).

\subsection{Financieras y seguros}

Esta agrupación reúne actividades financieras y de seguros. En la CIIU Rev. 4 A.C. corresponde a K; divisiones 64--66.

\begin{table}[H]
\centering
\caption{Financieras y seguros: PIB, ocupados, horas y productividad laboral, 2010--2025}
\label{tab:agg25_k_productividad}
\small
\begin{tabular}{lrrr}
\toprule
Indicador & 2010 & 2025 & Crec. anual \\
\midrule
PIB real  (Billones de pesos de 2015) & 22,3 & 53,1 & 5,95\% \\
Ocupados (Millones) & 0,2 & 0,4 & 4,49\% \\
PIB por trabajador (Millones de pesos de 2015) & 101,0 & 124,5 & 1,40\% \\
Horas semanales por trabajador & 45,2 & 44,2 & -0,15\% \\
PIB por hora trabajada (Miles de pesos de 2015) & 43,0 & 54,2 & 1,56\% \\
\bottomrule
\end{tabular}
\end{table}

\textbf{Frente al agregado nacional, la comparación variable por variable muestra diferencias relevantes.} El PIB real de la actividad registró una tasa anualizada de 5,95\%, por encima del agregado (3,16\%). La ocupación en la actividad registró una tasa anualizada de 4,49\%, por encima del agregado (2,09\%). El PIB por trabajador registró una tasa anualizada de 1,40\%, por encima del agregado (1,05\%). Las horas semanales por trabajador registraron una tasa anualizada de -0,15\%, una caída menos pronunciada que la del agregado (-0,45\%). El PIB por hora trabajada registró una tasa anualizada de 1,56\%, muy cerca del agregado (1,51\%).

\textbf{En niveles de 2025, el tamaño relativo y la productividad de la actividad económica también importan.} La actividad representó 5,2\% del PIB real agregado, con 53,1 billones de pesos de 2015, y concentró 1,9\% de los ocupados, con 0,4 millones de personas. El PIB por trabajador fue 124,5 millones de pesos de 2015 por ocupado, por encima del agregado (44,5 millones). Las horas semanales por trabajador fueron 44,2, muy cerca del agregado (43,7). El PIB por hora trabajada fue 54,2 mil pesos de 2015 por hora, por encima del agregado (19,6 mil).

\subsection{Inmobiliarias}

Esta agrupación reúne actividades inmobiliarias. En la CIIU Rev. 4 A.C. corresponde a L; división 68.

\begin{table}[H]
\centering
\caption{Inmobiliarias: PIB, ocupados, horas y productividad laboral, 2010--2025}
\label{tab:agg25_l_productividad}
\small
\begin{tabular}{lrrr}
\toprule
Indicador & 2010 & 2025 & Crec. anual \\
\midrule
PIB real  (Billones de pesos de 2015) & 59,9 & 90,1 & 2,75\% \\
Ocupados (Millones) & 0,2 & 0,3 & 2,64\% \\
PIB por trabajador (Millones de pesos de 2015) & 280,0 & 284,6 & 0,11\% \\
Horas semanales por trabajador & 51,1 & 45,8 & -0,73\% \\
PIB por hora trabajada (Miles de pesos de 2015) & 105,3 & 119,5 & 0,85\% \\
\bottomrule
\end{tabular}
\end{table}

\textbf{Frente al agregado nacional, la comparación variable por variable muestra diferencias relevantes.} El PIB real de la actividad registró una tasa anualizada de 2,75\%, por debajo del agregado (3,16\%). La ocupación en la actividad registró una tasa anualizada de 2,64\%, por encima del agregado (2,09\%). El PIB por trabajador registró una tasa anualizada de 0,11\%, por debajo del agregado (1,05\%). Las horas semanales por trabajador registraron una tasa anualizada de -0,73\%, una caída más pronunciada que la del agregado (-0,45\%). El PIB por hora trabajada registró una tasa anualizada de 0,85\%, por debajo del agregado (1,51\%).

\textbf{En niveles de 2025, el tamaño relativo y la productividad de la actividad económica también importan.} La actividad representó 8,8\% del PIB real agregado, con 90,1 billones de pesos de 2015, y concentró 1,4\% de los ocupados, con 0,3 millones de personas. El PIB por trabajador fue 284,6 millones de pesos de 2015 por ocupado, por encima del agregado (44,5 millones). Las horas semanales por trabajador fueron 45,8, muy cerca del agregado (43,7). El PIB por hora trabajada fue 119,5 mil pesos de 2015 por hora, por encima del agregado (19,6 mil).

\subsection{Profesionales y administrativas}

Esta agrupación reúne actividades profesionales, científicas y técnicas, y servicios administrativos y de apoyo. En la CIIU Rev. 4 A.C. corresponde a M y N; divisiones 69--82.

\begin{table}[H]
\centering
\caption{Profesionales y administrativas: PIB, ocupados, horas y productividad laboral, 2010--2025}
\label{tab:agg25_m_n_productividad}
\small
\begin{tabular}{lrrr}
\toprule
Indicador & 2010 & 2025 & Crec. anual \\
\midrule
PIB real  (Billones de pesos de 2015) & 45,4 & 70,1 & 2,94\% \\
Ocupados (Millones) & 0,8 & 1,7 & 4,94\% \\
PIB por trabajador (Millones de pesos de 2015) & 56,6 & 42,4 & -1,90\% \\
Horas semanales por trabajador & 41,9 & 39,6 & -0,38\% \\
PIB por hora trabajada (Miles de pesos de 2015) & 26,0 & 20,6 & -1,52\% \\
\bottomrule
\end{tabular}
\end{table}

\textbf{Frente al agregado nacional, la comparación variable por variable muestra diferencias relevantes.} El PIB real de la actividad registró una tasa anualizada de 2,94\%, por debajo del agregado (3,16\%). La ocupación en la actividad registró una tasa anualizada de 4,94\%, por encima del agregado (2,09\%). El PIB por trabajador registró una tasa anualizada de -1,90\%, por debajo del agregado (1,05\%). Las horas semanales por trabajador registraron una tasa anualizada de -0,38\%, muy cerca de la caída agregada (-0,45\%). El PIB por hora trabajada registró una tasa anualizada de -1,52\%, por debajo del agregado (1,51\%).

\textbf{En niveles de 2025, el tamaño relativo y la productividad de la actividad económica también importan.} La actividad representó 6,9\% del PIB real agregado, con 70,1 billones de pesos de 2015, y concentró 7,2\% de los ocupados, con 1,7 millones de personas. El PIB por trabajador fue 42,4 millones de pesos de 2015 por ocupado, muy cerca del agregado (44,5 millones). Las horas semanales por trabajador fueron 39,6, por debajo del agregado (43,7). El PIB por hora trabajada fue 20,6 mil pesos de 2015 por hora, por encima del agregado (19,6 mil).

\subsection{Administración pública}

Esta agrupación reúne administración pública y defensa, y planes de seguridad social de afiliación obligatoria. En la CIIU Rev. 4 A.C. corresponde a O; división 84.

\begin{table}[H]
\centering
\caption{Administración pública: PIB, ocupados, horas y productividad laboral, 2010--2025}
\label{tab:agg25_o_productividad}
\small
\begin{tabular}{lrrr}
\toprule
Indicador & 2010 & 2025 & Crec. anual \\
\midrule
PIB real  (Billones de pesos de 2015) & 31,3 & 62,7 & 4,75\% \\
Ocupados (Millones) & 0,5 & 0,8 & 3,30\% \\
PIB por trabajador (Millones de pesos de 2015) & 62,7 & 77,2 & 1,40\% \\
Horas semanales por trabajador & 50,9 & 46,2 & -0,65\% \\
PIB por hora trabajada (Miles de pesos de 2015) & 23,7 & 32,2 & 2,06\% \\
\bottomrule
\end{tabular}
\end{table}

\textbf{Frente al agregado nacional, la comparación variable por variable muestra diferencias relevantes.} El PIB real de la actividad registró una tasa anualizada de 4,75\%, por encima del agregado (3,16\%). La ocupación en la actividad registró una tasa anualizada de 3,30\%, por encima del agregado (2,09\%). El PIB por trabajador registró una tasa anualizada de 1,40\%, por encima del agregado (1,05\%). Las horas semanales por trabajador registraron una tasa anualizada de -0,65\%, una caída más pronunciada que la del agregado (-0,45\%). El PIB por hora trabajada registró una tasa anualizada de 2,06\%, por encima del agregado (1,51\%).

\textbf{En niveles de 2025, el tamaño relativo y la productividad de la actividad económica también importan.} La actividad representó 6,1\% del PIB real agregado, con 62,7 billones de pesos de 2015, y concentró 3,5\% de los ocupados, con 0,8 millones de personas. El PIB por trabajador fue 77,2 millones de pesos de 2015 por ocupado, por encima del agregado (44,5 millones). Las horas semanales por trabajador fueron 46,2, por encima del agregado (43,7). El PIB por hora trabajada fue 32,2 mil pesos de 2015 por hora, por encima del agregado (19,6 mil).

\subsection{Educación}

Esta agrupación reúne educación. En la CIIU Rev. 4 A.C. corresponde a P; división 85.

\begin{table}[H]
\centering
\caption{Educación: PIB, ocupados, horas y productividad laboral, 2010--2025}
\label{tab:agg25_p_productividad}
\small
\begin{tabular}{lrrr}
\toprule
Indicador & 2010 & 2025 & Crec. anual \\
\midrule
PIB real  (Billones de pesos de 2015) & 31,1 & 52,7 & 3,59\% \\
Ocupados (Millones) & 0,7 & 1,0 & 2,01\% \\
PIB por trabajador (Millones de pesos de 2015) & 41,6 & 52,3 & 1,54\% \\
Horas semanales por trabajador & 37,4 & 39,1 & 0,28\% \\
PIB por hora trabajada (Miles de pesos de 2015) & 21,4 & 25,8 & 1,26\% \\
\bottomrule
\end{tabular}
\end{table}

\textbf{Frente al agregado nacional, la comparación variable por variable muestra diferencias relevantes.} El PIB real de la actividad registró una tasa anualizada de 3,59\%, por encima del agregado (3,16\%). La ocupación en la actividad registró una tasa anualizada de 2,01\%, muy cerca del agregado (2,09\%). El PIB por trabajador registró una tasa anualizada de 1,54\%, por encima del agregado (1,05\%). Las horas semanales por trabajador registraron una tasa anualizada de 0,28\%, por encima del agregado (-0,45\%). El PIB por hora trabajada registró una tasa anualizada de 1,26\%, por debajo del agregado (1,51\%).

\textbf{En niveles de 2025, el tamaño relativo y la productividad de la actividad económica también importan.} La actividad representó 5,2\% del PIB real agregado, con 52,7 billones de pesos de 2015, y concentró 4,4\% de los ocupados, con 1,0 millones de personas. El PIB por trabajador fue 52,3 millones de pesos de 2015 por ocupado, por encima del agregado (44,5 millones). Las horas semanales por trabajador fueron 39,1, por debajo del agregado (43,7). El PIB por hora trabajada fue 25,8 mil pesos de 2015 por hora, por encima del agregado (19,6 mil).

\subsection{Salud y servicios sociales}

Esta agrupación reúne actividades de atención de la salud humana y de servicios sociales. En la CIIU Rev. 4 A.C. corresponde a Q; divisiones 86--88.

\begin{table}[H]
\centering
\caption{Salud y servicios sociales: PIB, ocupados, horas y productividad laboral, 2010--2025}
\label{tab:agg25_q_productividad}
\small
\begin{tabular}{lrrr}
\toprule
Indicador & 2010 & 2025 & Crec. anual \\
\midrule
PIB real  (Billones de pesos de 2015) & 23,6 & 49,8 & 5,11\% \\
Ocupados (Millones) & 0,7 & 1,0 & 2,40\% \\
PIB por trabajador (Millones de pesos de 2015) & 33,2 & 49,1 & 2,65\% \\
Horas semanales por trabajador & 44,0 & 44,5 & 0,06\% \\
PIB por hora trabajada (Miles de pesos de 2015) & 14,5 & 21,3 & 2,58\% \\
\bottomrule
\end{tabular}
\end{table}

\textbf{Frente al agregado nacional, la comparación variable por variable muestra diferencias relevantes.} El PIB real de la actividad registró una tasa anualizada de 5,11\%, por encima del agregado (3,16\%). La ocupación en la actividad registró una tasa anualizada de 2,40\%, por encima del agregado (2,09\%). El PIB por trabajador registró una tasa anualizada de 2,65\%, por encima del agregado (1,05\%). Las horas semanales por trabajador registraron una tasa anualizada de 0,06\%, por encima del agregado (-0,45\%). El PIB por hora trabajada registró una tasa anualizada de 2,58\%, por encima del agregado (1,51\%).

\textbf{En niveles de 2025, el tamaño relativo y la productividad de la actividad económica también importan.} La actividad representó 4,9\% del PIB real agregado, con 49,8 billones de pesos de 2015, y concentró 4,4\% de los ocupados, con 1,0 millones de personas. El PIB por trabajador fue 49,1 millones de pesos de 2015 por ocupado, por encima del agregado (44,5 millones). Las horas semanales por trabajador fueron 44,5, muy cerca del agregado (43,7). El PIB por hora trabajada fue 21,3 mil pesos de 2015 por hora, por encima del agregado (19,6 mil).

\subsection{Artes, entretenimiento y otros servicios}

Esta agrupación reúne actividades artísticas, entretenimiento, recreación y otros servicios. En la CIIU Rev. 4 A.C. corresponde a R y S; divisiones 90--96.

\begin{table}[H]
\centering
\caption{Artes, entretenimiento y otros servicios: PIB, ocupados, horas y productividad laboral, 2010--2025}
\label{tab:agg25_r_s_productividad}
\small
\begin{tabular}{lrrr}
\toprule
Indicador & 2010 & 2025 & Crec. anual \\
\midrule
PIB real  (Billones de pesos de 2015) & 11,0 & 40,6 & 9,09\% \\
Ocupados (Millones) & 0,8 & 1,2 & 2,97\% \\
PIB por trabajador (Millones de pesos de 2015) & 13,8 & 32,8 & 5,95\% \\
Horas semanales por trabajador & 39,2 & 39,6 & 0,06\% \\
PIB por hora trabajada (Miles de pesos de 2015) & 6,8 & 16,0 & 5,89\% \\
\bottomrule
\end{tabular}
\end{table}

\textbf{Frente al agregado nacional, la comparación variable por variable muestra diferencias relevantes.} El PIB real de la actividad registró una tasa anualizada de 9,09\%, por encima del agregado (3,16\%). La ocupación en la actividad registró una tasa anualizada de 2,97\%, por encima del agregado (2,09\%). El PIB por trabajador registró una tasa anualizada de 5,95\%, por encima del agregado (1,05\%). Las horas semanales por trabajador registraron una tasa anualizada de 0,06\%, por encima del agregado (-0,45\%). El PIB por hora trabajada registró una tasa anualizada de 5,89\%, por encima del agregado (1,51\%).

\textbf{En niveles de 2025, el tamaño relativo y la productividad de la actividad económica también importan.} La actividad representó 4,0\% del PIB real agregado, con 40,6 billones de pesos de 2015, y concentró 5,4\% de los ocupados, con 1,2 millones de personas. El PIB por trabajador fue 32,8 millones de pesos de 2015 por ocupado, por debajo del agregado (44,5 millones). Las horas semanales por trabajador fueron 39,6, por debajo del agregado (43,7). El PIB por hora trabajada fue 16,0 mil pesos de 2015 por hora, por debajo del agregado (19,6 mil).

\subsection{Hogares como empleadores}

Esta agrupación reúne actividades de los hogares como empleadores y producción no diferenciada de los hogares para uso propio. En la CIIU Rev. 4 A.C. corresponde a T; divisiones 97--98.

\begin{table}[H]
\centering
\caption{Hogares como empleadores: PIB, ocupados, horas y productividad laboral, 2010--2025}
\label{tab:agg25_t_productividad}
\small
\begin{tabular}{lrrr}
\toprule
Indicador & 2010 & 2025 & Crec. anual \\
\midrule
PIB real  (Billones de pesos de 2015) & 4,3 & 6,7 & 3,06\% \\
Ocupados (Millones) & 0,7 & 0,7 & 0,42\% \\
PIB por trabajador (Millones de pesos de 2015) & 6,4 & 9,4 & 2,63\% \\
Horas semanales por trabajador & 45,9 & 40,5 & -0,83\% \\
PIB por hora trabajada (Miles de pesos de 2015) & 2,7 & 4,5 & 3,49\% \\
\bottomrule
\end{tabular}
\end{table}

\textbf{Frente al agregado nacional, la comparación variable por variable muestra diferencias relevantes.} El PIB real de la actividad registró una tasa anualizada de 3,06\%, muy cerca del agregado (3,16\%). La ocupación en la actividad registró una tasa anualizada de 0,42\%, por debajo del agregado (2,09\%). El PIB por trabajador registró una tasa anualizada de 2,63\%, por encima del agregado (1,05\%). Las horas semanales por trabajador registraron una tasa anualizada de -0,83\%, una caída más pronunciada que la del agregado (-0,45\%). El PIB por hora trabajada registró una tasa anualizada de 3,49\%, por encima del agregado (1,51\%).

\textbf{En niveles de 2025, el tamaño relativo y la productividad de la actividad económica también importan.} La actividad representó 0,7\% del PIB real agregado, con 6,7 billones de pesos de 2015, y concentró 3,1\% de los ocupados, con 0,7 millones de personas. El PIB por trabajador fue 9,4 millones de pesos de 2015 por ocupado, por debajo del agregado (44,5 millones). Las horas semanales por trabajador fueron 40,5, por debajo del agregado (43,7). El PIB por hora trabajada fue 4,5 mil pesos de 2015 por hora, por debajo del agregado (19,6 mil).

{\footnotesize Nota general: para facilitar la lectura, el informe usa PIB por actividad económica; en sentido estricto, el numerador corresponde al valor agregado bruto de cada actividad reportado por el DANE. Ocupados expandidos con el factor \texttt{fex}. Las horas semanales promedio se calculan como horas anuales totales divididas por ocupados y por 52 semanas. Se excluye 2020 por no contar con GEIH anual comparable en la base del proyecto. Fuente: cálculos propios con DANE y GEIH.}
